\section{User interface design}

\subsection{Design direction}
Platform hiện thực các use case đã phân tích trong tài liệu này bằng một hệ
thống fullstack nhẹ, có profile local để demo nhanh và target triển khai GCP.
Front-end sử dụng HTML, CSS và JavaScript thuần; back-end sử dụng Node.js. Khi
deploy lên GCP, service Node.js được đặt trong VM/container và file SQLite/WAL nằm
trên persistent disk; khi chạy local, cùng SQLite schema được dùng để nhóm kiểm
thử end-to-end trên một máy nhưng vẫn thể hiện rõ luồng dữ liệu, trạng thái xe,
yêu cầu thuê, phiếu thuê và báo cáo vận hành.

Giao diện được thiết kế theo phong cách Civic Mobility Command Center: rail điều
hướng tối có thể cuộn tới từng vùng chức năng, workspace sáng, panel/form/table
sắc cạnh và mật độ thông tin phù hợp phần mềm điều phối vận hành. Bảng màu vẫn
bám chủ đề đô thị xanh Ecopark nhưng
không phụ thuộc vào một mảng xanh đơn điệu; các accent xanh nước và amber được
dùng để phân biệt trạng thái, bản đồ, loại xe và hành động chính. Lớp giao diện
được tinh chỉnh với thanh trên sticky dạng floating surface, session chip, menu
dropdown custom có trạng thái option rõ, trạng thái focus dễ nhận biết và bố cục
mobile tự trở về đầu trang khi chuyển màn hình làm việc. Thanh trên dùng nền gần
đặc, khoảng hở 14px và lớp mờ quanh mép để tách khỏi nội dung bên dưới mà không
làm rối khi cuộn hoặc mở dropdown. Giao diện cũng hỗ trợ ba chế độ \textit{Sáng}, \textit{Tối} và
\textit{Hệ thống}; theme được lưu ở trình duyệt, triển khai bằng design tokens
cho panel, bảng, dropdown và chart; bản đồ Leaflet có lớp xử lý màu riêng cho
dark mode còn scene Three.js giữ đủ tương phản để không mất khả năng quan sát khi demo. Trên mobile,
giao diện bỏ topbar riêng,
đưa scene 3D lên đầu hero, nén các metric thành một hàng ngang và rút rail điều
hướng tối thành dock icon-only nổi ở đáy màn hình, bao gồm cả refresh/theme/logout
để không cần một thanh topbar thứ hai. Ở desktop rộng, các panel có bảng dài như
\textit{Báo cáo \& CSV}, \textit{Yêu cầu nhận xe}, \textit{Danh sách trả xe} và
\textit{Đội xe} chiếm toàn bộ bề ngang workspace để chart và bảng không bị cắt.
Các dropdown ở bộ lọc báo cáo, nhận xe, trả xe, bãi trả và trạng thái đội xe
không dùng dropdown native của trình duyệt mà dùng cùng component custom có icon,
mô tả và trạng thái chọn; khi mở trong bảng, panel được nâng lớp z-index để menu
không bị cắt bởi panel kế tiếp. Chart tổng thu/phụ thu dùng cùng thang tiền với
chart doanh thu để so sánh tuần/tháng không bị lệch nghĩa và các thanh trạng thái
được bo tròn đồng nhất thay vì để phần màu bên trong thành hình chữ nhật.
Ở các viewport vận hành hẹp hoặc tablet/desktop nhỏ, hero staff/admin chuyển
sang một cột, dashboard chart giảm số cột và các bảng vận hành được chuyển thành
row-card có nhãn từng trường, nên báo cáo, yêu cầu nhận xe và đội xe vẫn đọc được
mà không cần kéo ngang trong panel. Phần đầu trang có một scene 3D nhỏ bằng
Three.js theo phong cách low-poly, mô phỏng Bike Hub isometric với đường nội khu,
bike lane, canopy, dock/rack, đèn trạng thái, hồ nước, bờ cỏ, dải cây xanh và ba
kiểu xe khác nhau. Bản scene mới mở rộng mặt hồ, bờ kè, cụm cây lớn quanh hồ và
hàng cây ven đường để phần minh hoạ có chiều sâu hơn trong khung hero hẹp. Scene
có thêm xe chạy vòng trên hành lang đường, người đi
bộ/kiểm tra xe/em nhỏ vẫy tay và sway cây nhẹ để tạo loop thị giác mượt. Scene
này dùng camera/frustum responsive để tự fit trong khung đăng nhập, hero dashboard
và mobile, thay cho cách phóng canvas bằng CSS; khi render lại cùng màn hình, app
giữ node canvas cũ để tránh nhấp nháy. Scene được đóng thành một dải preview cân
với cột nội dung ở màn hình đăng nhập và chỉ đóng vai trò nhấn thị giác, không
thay thế các bảng, form và trạng thái nghiệp vụ chính.

Các thẻ xe trong luồng khách hàng dùng minh hoạ SVG theo loại xe thay vì dùng
một icon chung cho mọi dữ liệu. City bike, Family tandem và Child-seat bike có
màu nhấn và hình dáng khác nhau để người dùng nhận biết loại xe nhanh hơn trước
khi gửi yêu cầu thuê. Card xe dùng bố cục ngang: minh họa loại xe ở cột trái, mã
xe, loại xe, trạng thái sẵn sàng và nút thuê gom ở cột phải để giảm khoảng trống
trong danh sách lựa chọn. Trên màn hình hẹp, hàng chọn thời lượng và nút thuê
được đưa xuống dùng toàn bộ bề ngang card để tránh bó chữ hoặc làm lệch viewport.
Khi xe không thuê được, card hiển thị rõ nguyên nhân như \textit{Chờ bạn nhận},
\textit{Bạn đang thuê}, \textit{Người khác giữ} hoặc \textit{Người khác thuê}
thay vì chỉ ghi chung là đang thuê.

Lớp chuyển động sử dụng GSAP cục bộ từ dependency npm. Các hiệu ứng vào trang được
gom bằng timeline, ưu tiên transform và opacity để giảm chi phí layout; các tương
tác nhỏ khi hover và toast cũng dùng GSAP với overwrite phù hợp. Giao diện tôn trọng
\texttt{prefers-reduced-motion}: khi người dùng giảm chuyển động, các hiệu ứng
được bỏ qua hoặc rút về trạng thái tĩnh. Các chi tiết giao diện đã được rà lại qua
các lượt chụp màn hình và tinh chỉnh, bao gồm bố cục màn hình đăng nhập, kích thước
hero trên mobile, độ dày bảng quản trị, vị trí toast, vị trí cuộn sau đăng nhập và cân
bằng giữa scene 3D với nội dung nghiệp vụ.

Màn hình khách hàng sử dụng bản đồ thật bằng Leaflet/OpenStreetMap để hiển thị
vị trí người dùng demo và marker các bãi xe Ecopark. Khách có thể dùng GPS demo
hoặc chuyển sang vị trí nhập tay, chọn phạm vi tìm kiếm, lọc loại xe bằng custom
dropdown HTML/CSS, chọn marker hoặc thẻ bãi để cập nhật danh sách xe sẵn sàng
tại bãi đó. Dropdown loại khách trong form đăng ký cũng dùng control custom thay
cho popup native của trình duyệt. Dropdown thời lượng trong card xe cũng dùng
cùng mẫu custom control, chỉ hiển thị nhãn giá ngắn trong trigger và mở menu có
mô tả đầy đủ khi cần chọn 1/2/3 giờ. Các dropdown vẫn giữ chiều rộng ổn định trên
desktop/mobile, không tạo scroll ngang và có trạng thái option đang chọn rõ ràng.
Khi không có bãi trong phạm vi đã chọn, giao diện hiển thị trạng thái rỗng rõ ràng để kích hoạt
alternative flow tìm kiếm thủ công/mở rộng phạm vi. Nếu tile bản đồ
không tải được trong môi trường offline, khung fallback nội bộ vẫn giữ pin bãi xe
để luồng nghiệp vụ không bị gián đoạn.

Workspace khách hàng gom vị trí tìm kiếm, bản đồ, bãi gần nhất, danh sách xe rảnh
và nút gửi yêu cầu vào cùng vùng \textit{Thuê xe}. Các vùng \textit{Lượt thuê}
và \textit{Tài khoản} được tách xuống dưới/ra rail điều hướng để màn đầu không
giống bảng quản trị thu nhỏ. Trạng thái GPS hoặc vị trí nhập tay được hiển thị
như chip overlay trong bản đồ, vì đây là ngữ cảnh của bản đồ thay vì một cảnh báo
riêng bên ngoài bộ lọc.

\subsection{Interface evidence}
Các hình dưới đây là screenshot trực tiếp từ platform hiện tại, dùng để chứng
minh UI đã có màn đăng nhập/đăng ký, workspace khách hàng trên desktop/mobile,
workspace vận hành và màn GPS hỗ trợ demo use case.

\begin{figure}[H]
\centering
\includegraphics[width=0.92\textwidth,height=0.40\textheight,keepaspectratio]{auth_desktop.png}
\caption{Màn đăng nhập/đăng ký và 3D Bike Hub preview}
\end{figure}

\subsection{UC001 screen evidence}
Chuỗi ảnh dưới đây được chụp trực tiếp bằng script
\texttt{npm run screenshots:uc001}. Chuỗi này minh chứng basic flow UC001 và
alternative flow 1: hệ thống chặn form đăng ký khi thiếu thông tin bắt buộc,
hiển thị trường lỗi và giữ khách ở bước nhập liệu.

\begin{figure}[H]
\centering
\includegraphics[width=0.48\textwidth,height=0.52\textheight,keepaspectratio]{uc001_flow/uc001_01_register_form.png}
\caption{UC001 bước 3: khách mở form tạo tài khoản với họ tên, email, số điện thoại, mật khẩu, CCCD/CMND, địa chỉ và loại khách}
\end{figure}

\begin{figure}[H]
\centering
\includegraphics[width=0.92\textwidth,height=0.44\textheight,keepaspectratio]{uc001_flow/uc001_02_missing_required_fields.png}
\caption{UC001 alternative flow 1: thiếu email bắt buộc, UI giữ form tại bước nhập và hiển thị lỗi rõ ràng}
\end{figure}

\begin{figure}[H]
\centering
\begin{minipage}{0.31\textwidth}
\centering
\includegraphics[width=\linewidth,height=0.42\textheight,keepaspectratio]{uc001_flow/uc001_03_account_created_unverified.png}\\
{\scriptsize Tài khoản tạo xong, email chưa xác minh}
\end{minipage}\hfill
\begin{minipage}{0.31\textwidth}
\centering
\includegraphics[width=\linewidth,height=0.42\textheight,keepaspectratio]{uc001_flow/uc001_04_email_code_ready.png}\\
{\scriptsize Mã xác minh demo-local được tạo trong platform}
\end{minipage}\hfill
\begin{minipage}{0.31\textwidth}
\centering
\includegraphics[width=\linewidth,height=0.42\textheight,keepaspectratio]{uc001_flow/uc001_05_email_verified.png}\\
{\scriptsize Email chuyển sang trạng thái đã xác minh}
\end{minipage}
\caption{UC001 bước 5--10: trạng thái tài khoản, mã xác minh email và hồ sơ đủ điều kiện thuê}
\end{figure}

\begin{figure}[H]
\centering
\includegraphics[width=0.92\textwidth,height=0.42\textheight,keepaspectratio]{customer_rent_desktop.png}
\caption{Workspace khách hàng: tìm bãi, bản đồ, chọn xe và gửi yêu cầu thuê}
\end{figure}

\clearpage
\begin{figure}[p]
\centering
\begin{minipage}{0.31\textwidth}
\centering
\includegraphics[width=\linewidth,height=0.58\textheight,keepaspectratio]{customer_mobile_overview.png}\\
{\scriptsize Tổng quan}
\end{minipage}\hfill
\begin{minipage}{0.31\textwidth}
\centering
\includegraphics[width=\linewidth,height=0.58\textheight,keepaspectratio]{customer_mobile_search.png}\\
{\scriptsize Tìm bãi}
\end{minipage}\hfill
\begin{minipage}{0.31\textwidth}
\centering
\includegraphics[width=\linewidth,height=0.58\textheight,keepaspectratio]{customer_mobile_bikes.png}\\
{\scriptsize Chọn xe}
\end{minipage}
\caption{UI điện thoại: 3D hero, metric compact, dock icon-only và luồng thuê xe}
\end{figure}
\clearpage

\begin{figure}[p]
\centering
\begin{minipage}{0.96\textwidth}
\centering
\includegraphics[width=\linewidth,height=0.32\textheight,keepaspectratio]{operations_overview.png}\\[2mm]
\raggedright{\small Tổng quan ca vận hành, yêu cầu nhận xe và pipeline trả xe}
\end{minipage}\\[4mm]
\begin{minipage}{0.96\textwidth}
\centering
\includegraphics[width=\linewidth,height=0.32\textheight,keepaspectratio]{operations_report.png}\\[2mm]
\raggedright{\small Dashboard báo cáo, trạng thái đội xe và bộ lọc xuất CSV}
\end{minipage}
\caption{Workspace vận hành: nhận xe, trả xe, quản lý đội xe và báo cáo}
\end{figure}
\clearpage

\begin{figure}[H]
\centering
\includegraphics[width=0.86\textwidth,height=0.40\textheight,keepaspectratio]{gps_demo.png}
\caption{Route \texttt{/gd}: điều khiển GPS người dùng và đồng hồ hệ thống cho demo UC002/UC004}
\end{figure}

\subsection{UI screen-flow diagram}
\begin{figure}[H]
\centering
\resizebox{\textwidth}{!}{%
\begin{tikzpicture}[
  lane/.style={font=\footnotesize\bfseries,align=right,text=Ink},
  screen/.style={draw=EcoparkGreen,thick,rounded corners=2pt,fill=SoftGreen,minimum width=3.0cm,minimum height=0.82cm,align=center,font=\small,inner xsep=5pt},
  ops/.style={draw=EcoparkAmber,thick,rounded corners=2pt,fill=yellow!10,minimum width=3.0cm,minimum height=0.82cm,align=center,font=\small,inner xsep=5pt},
  aid/.style={draw=EcoparkTeal,thick,rounded corners=2pt,fill=cyan!8,minimum width=3.0cm,minimum height=0.82cm,align=center,font=\small,inner xsep=5pt},
  arr/.style={-{Latex[length=2.2mm]},draw=Ink,thick},
  support/.style={-{Latex[length=2.0mm]},draw=EcoparkTeal,thick,dashed}
]
\node[lane] at (0,1.60) {Customer};
\node[lane] at (0,0.15) {Staff/Admin};
\node[lane] at (0,-1.30) {Demo};

\node[screen] (authCustomer) at (2.35,1.60) {Login\\Register};
\node[screen] (rent) at (5.75,1.60) {Tìm bãi\\Chọn xe};
\node[screen] (history) at (9.15,1.60) {Lượt thuê\\Tài khoản};
\node[ops] (authStaff) at (2.35,0.15) {Login\\Register};
\node[ops] (pickup) at (5.75,0.15) {Nhận xe\\Giữ cọc};
\node[ops] (return) at (9.15,0.15) {Trả xe\\Xuất vé};
\node[ops] (manage) at (12.55,0.15) {Báo cáo\\Đội xe/Audit};
\node[aid] (director) at (2.35,-1.30) {\texttt{/gd}\\Demo director};
\node[aid] (gps) at (5.75,-1.30) {GPS người dùng\\theo đường};
\node[aid] (clock) at (9.15,-1.30) {Giờ hệ thống\\phụ thu};

\draw[arr] (authCustomer) -- (rent);
\draw[arr] (rent) -- (history);
\draw[arr] (authStaff) -- (pickup);
\draw[arr] (pickup) -- (return);
\draw[arr] (return) -- (manage);
\draw[support] (director) -- (gps);
\draw[support] (gps) -- (clock);
\end{tikzpicture}
}
\caption{Screen-flow chính của platform và route hỗ trợ demo GPS}
\end{figure}

Form đăng ký khách có gợi ý nhập liệu cho họ tên, số điện thoại, CCCD/CMND, địa
chỉ và nút hiện/ẩn mật khẩu để người dùng tự kiểm tra trước khi gửi. API đăng ký kiểm tra họ tên không chứa số/ký tự lạ, số điện thoại Việt Nam
dạng \texttt{0xxxxxxxxx} hoặc \texttt{+84xxxxxxxxx}, CCCD/CMND gồm 9 hoặc 12 chữ
số, đồng thời chặn số điện thoại hoặc giấy tờ định danh đã được dùng bởi tài
khoản khác. Theo UC001, CCCD/CMND hợp lệ là điều kiện tối thiểu trước khi thuê;
staff vẫn đối chiếu giấy tờ tại bước UC003 handover, còn trạng thái thẻ cư dân
pending/rejected chỉ ảnh hưởng ưu đãi cư dân chứ không chặn quyền thuê như khách thường.

Màn hình vận hành có pipeline trả xe riêng, luôn hiển thị bốn bước: nhận xe,
khách xác nhận bãi trả bằng GPS, kiểm tra tình trạng xe và xuất vé. Bảng chi
tiết bên dưới vẫn xử lý các input nghiệp vụ đầy đủ như bãi trả, giờ trả, tình
trạng xe và ghi chú kiểm tra, nhưng nút xuất vé chỉ mở khi customer đã xác nhận
vị trí trong bán kính bãi trả, giúp actor thấy rõ luồng UC004 ngay cả khi chưa
có rental active.

Để hỗ trợ trình bày use case trước lớp, ứng dụng có thêm route demo
\texttt{/gd}. Đây là demo director, không phải nghiệp vụ production mới: người
trình bày tìm/chọn từng tài khoản khách, chọn phiên khách đang chờ nhận xe hoặc
đang thuê xe khi có, chọn bãi đích, kéo thả marker GPS người dùng hoặc bấm snap
gần bãi. Hệ thống sẽ đưa marker về điểm gần nhất trên tuyến đường demo nội bộ và
hiển thị polyline đường đi tới bãi nhận/trả. Tọa độ sau khi snap hoặc kéo được
lưu theo từng customer qua API demo dùng chung. Khi marker \texttt{/gd} đang
chạy animation trên polyline đường, client stream các điểm trung gian theo
throttle, vì vậy tab khách hàng tương ứng đang chọn \texttt{GPS hiện tại} poll
cùng API và animate marker ``Bạn'', khoảng cách bãi và danh sách xe theo vị trí
đó thay vì chỉ nhảy tới điểm cuối. Khi chưa có lượt
active ở mode trả xe, UI vẫn cho kéo user đã chọn để chuẩn bị demo và hiển thị
empty state cho phần phiên UC. Cùng màn này có đồng hồ hệ thống nội bộ:
staff/admin có thể tua +15/+30/+60 phút hoặc reset về giờ thật để diễn trường
hợp quá hạn và phụ thu 30k/30 phút. Nhờ đó phần diễn UC002/UC004 không tạo cảm
giác xe bay thẳng qua hồ hoặc nhà, đồng thời phí phạt vẫn được tính bởi backend
theo cùng mốc giờ mà customer và staff nhìn thấy. Khi mở \texttt{/gd} trước lúc
danh sách tài khoản, bãi và phiên GPS sẵn sàng, giao diện hiển thị skeleton loading để không xuất hiện
một grid trắng trong lúc demo.

\subsection{Use-case to screen mapping}
\begin{longtable}{P{0.18\textwidth}P{0.34\textwidth}P{0.38\textwidth}}
\toprule
\textbf{Use case} & \textbf{Màn hình / chức năng web} & \textbf{Kết quả triển khai} \\
\midrule
UC001 & Đăng ký, đăng nhập, xác minh email, reset mật khẩu, hồ sơ khách, thẻ cư dân & Khách có tài khoản, CCCD/email/số điện thoại/địa chỉ, email đã xác minh, mật khẩu theo policy và trạng thái giảm giá cư dân nếu có thẻ hợp lệ. \\
UC002 & Danh sách bãi theo GPS demo/vị trí đã lưu/vị trí nhập tay, lọc loại xe, chọn xe, hủy request & Hệ thống hiển thị bãi trong phạm vi, số xe rảnh, trạng thái từng xe, tạo yêu cầu thuê với thời lượng 1/2/3 giờ ở mọi vị trí khách gửi và cho hủy request trước khi nhận xe. Vị trí hiện tại chỉ trở thành điều kiện bắt buộc ở bước handover. \\
UC003 & Bảng yêu cầu nhận xe cho Staff/Admin & Nhân sự được scope theo bãi, đối chiếu loại/số giấy tờ, nhận cọc 200k, ghi nhận giấy tờ giữ tại bãi, đổi sang xe rảnh cùng bãi khi cần, kiểm tra lại geofence khách ở bước handover, hủy yêu cầu không đủ điều kiện, tạo rental và chuyển xe sang trạng thái đang thuê; khi thao tác bị chặn, UI hiển thị lý do ngay trong panel nhận xe. \\
UC004 & Pipeline trả xe, bảng xe đang thuê và thao tác trả xe & Staff nhận xe tại bãi được phân công; Admin/Operator xử lý toàn hệ thống và hỗ trợ trả khác bãi. Customer phải xác nhận bãi trả bằng GPS trong bán kính bãi trước khi staff xuất vé; hệ thống tính phí thuê, chỉ giảm giá cư dân khi tài khoản là resident và thẻ cư dân verified, phụ thu trễ và hiển thị phiếu vừa xuất; nếu chưa đủ điều kiện xuất vé, UI hiện lý do thay vì chỉ khóa nút. \\
UC005 & Fleet table, quản lý bãi/xe, dashboard báo cáo, audit log & Admin/Operator quản lý bãi xe, xe, trạng thái xe, tra cứu vị trí xe, xem chart doanh thu/trạng thái đội xe/tổng thu và phụ thu, lọc thống kê bằng custom dropdown theo ngày/tuần/tháng, theo bãi/xe, xuất CSV và xem nhật ký đổi trạng thái. \\
\bottomrule
\end{longtable}

\section{Component implementation}

\subsection{Implemented modules}
\begin{longtable}{P{0.22\textwidth}P{0.28\textwidth}P{0.40\textwidth}}
\toprule
\textbf{Module} & \textbf{File chính} & \textbf{Vai trò triển khai} \\
\midrule
Frontend shell & \texttt{public/app.js}, \texttt{public/styles.css} & Render auth, customer, operations và \texttt{/gd}; quản lý state client, custom dropdown, map, table, form, report charts, audit timeline, demo clock, CSV download, inline operation errors và toast overlay không bị topbar che. \\
3D scene & \texttt{public/scene.js} & Mount scene Three.js low-poly, tự cleanup scene cũ, dùng frustum responsive và giữ \texttt{preserveDrawingBuffer} cho smoke test pixel. \\
API server & \texttt{server/app.js} & HTTP router, session cookie, role guard, static/vendor serving và các endpoint nghiệp vụ. \\
Database & \texttt{server/db.js} & SQLite schema, seed data, demo accounts, foreign keys và WAL mode; khi deploy GCP, file DB được đặt trên persistent disk. \\
Pricing & \texttt{server/pricing.js}, \texttt{server/app.js} & Tính phí cơ bản, phụ thu trễ và giảm giá cư dân theo policy đề bài; guard ticket-time yêu cầu customer\_type resident, cờ discount hợp lệ và resident card verified. \\
Backend tests & \texttt{test/app.test.js} & Test đăng ký, availability, resident discount guard, tạo request ngoài bán kính nhưng chặn handover theo GPS hiện tại, geofence return, handover, return ticket, reports và concurrent sessions. \\
Browser smoke tests & \texttt{scripts/ui-smoke.js}, \texttt{scripts/uc-flow-smoke.js} & Kiểm tra render UI, map/scene, overflow, console error, tab customer nhận marker GPS stream từ \texttt{/gd} và pipeline UC bằng Playwright. \\
\bottomrule
\end{longtable}

\subsection{Use-case implementation traceability}
\begin{longtable}{P{0.14\textwidth}P{0.30\textwidth}P{0.46\textwidth}}
\toprule
\textbf{UC} & \textbf{Component chính} & \textbf{Implementation note} \\
\midrule
UC001 & Auth UI, \texttt{registerCustomer}, \texttt{verifyEmail}, \texttt{resetPassword}, \texttt{submitResidentCard} & Custom customer-type dropdown, validation họ tên/phone/CCCD/password, PBKDF2 salted hash, email verification code, reset code, resident card verified/pending/rejected và blocked identity policy. \\
UC002 & Customer rental UI, station/bike APIs & Tìm bãi theo tọa độ hoặc vị trí đã lưu, chọn range/type, hiển thị held bike và tạo request pending pickup sau khi email/identity hợp lệ; API lưu tọa độ request nhưng chưa chặn theo bán kính bãi ở bước này. \\
UC003 & Pending request table, \texttt{handoverBike}, \texttt{exchangeRentalRequestBike} & Staff/Admin xác nhận giấy tờ, cọc 200k, scope request theo bãi nhân sự, kiểm tra GPS hiện tại của customer trước handover, đổi xe trước handover, tạo rental/deposit và đổi xe sang rented trong transaction; lỗi sai cọc/giấy tờ/bãi/bán kính được map thành thông báo tiếng Việt trong panel. \\
UC004 & Return pipeline, \texttt{confirmRentalReturn}, \texttt{returnBike}, pricing module & Customer xác nhận bãi trả trong bán kính bãi, staff nhận xe khác bãi khi đúng scope, tính ticket với resident discount guard tại thời điểm xuất vé, trả cọc, cập nhật bike station/status và ghi status log; nút xuất vé vẫn cho bấm khi chưa sẵn sàng để hiện lý do cần khách xác nhận. \\
UC005 & Fleet/report/audit panels, admin APIs & Quản lý bãi/xe, chặn sức chứa nhỏ hơn số xe, chặn đổi bãi xe đang thuê, lọc báo cáo bằng custom dropdown, hiển thị chart dùng cùng scale tiền, đọc audit log và export CSV. \\
\bottomrule
\end{longtable}

\section{API design and implementation}
Các API chính gồm nhóm xác thực, tra cứu bãi/xe, tạo yêu cầu thuê, giao xe, trả xe
và báo cáo. Dữ liệu được lưu trong SQLite/WAL theo các nhóm quan hệ đã nêu ở phần
\textit{Database design}; cùng DB này được dùng cho local demo và GCP deployment:
\begin{itemize}
  \item tài khoản và cư dân: USERS, CUSTOMER\_PROFILES, RESIDENT\_CARDS;
  \item bãi xe và đội xe: STATIONS, BIKE\_TYPES, BIKES;
  \item yêu cầu và lượt thuê: RENTAL\_REQUESTS, RENTALS;
  \item đặt cọc và vé thuê: RENTAL\_DEPOSITS, RENTAL\_TICKETS;
  \item truy vết vận hành và demo: BIKE\_STATUS\_LOGS, APP\_SETTINGS.
\end{itemize}

\subsection{API contract summary}
\begingroup
\small
\setlength{\tabcolsep}{4pt}
\begin{longtable}{P{0.12\textwidth}P{0.35\textwidth}P{0.16\textwidth}P{0.29\textwidth}}
\toprule
\textbf{Method} & \textbf{Endpoint} & \textbf{Role} & \textbf{Purpose} \\
\midrule
GET & \url{/api/session} & Public/session & Lấy user hiện tại từ session cookie. \\
POST & \url{/api/auth/register} & Public & Tạo tài khoản customer và hồ sơ định danh. \\
POST & \url{/api/auth/login} & Public & Đăng nhập và tạo session cookie. \\
POST & \url{/api/auth/logout} & Any session & Xóa session cookie hiện tại. \\
POST & \url{/api/auth/request-email-verification} & Customer & Tạo mã xác minh email demo-local, lưu dạng hash và thời hạn trong DB. \\
POST & \url{/api/auth/verify-email} & Customer & Xác minh email bằng mã demo-local trước khi gửi yêu cầu thuê. \\
POST & \url{/api/auth/request-password-reset} & Public & Tạo mã reset mật khẩu demo-local cho email active. \\
POST & \url{/api/auth/reset-password} & Public & Đặt lại mật khẩu theo policy và lưu hash PBKDF2 mới. \\
GET & \url{/api/bike-types} & Public & Lấy danh sách loại xe active. \\
GET & \url{/api/stations} & Public/customer & Lấy bãi xe, tồn kho tổng hợp và khoảng cách theo tọa độ. \\
GET & \url{/api/stations/:id/bikes} & Public/customer & Lấy xe trong bãi, trạng thái giữ chỗ và available. \\
GET & \url{/api/demo-clock} & Customer, Staff, Admin & Lấy giờ hệ thống nội bộ và mức phụ thu trễ cho demo. \\
POST & \url{/api/demo-clock} & Staff/Admin & Tua hoặc reset đồng hồ nội bộ dùng cho request, handover, return và audit log. \\
POST & \url{/api/rental-requests} & Customer & Tạo yêu cầu thuê xe pending pickup ở mọi vị trí khách gửi; vị trí được lưu để đối chiếu/fallback, còn điều kiện gần bãi được guard ở handover. \\
POST & \url{/api/rental-requests/:id/cancel} & Customer, Staff, Admin & Hủy request pending và giải phóng xe. \\
POST & \url{/api/staff/rental-requests/:id/exchange-bike} & Staff/Admin & Đổi request pending sang xe available cùng bãi trước handover. \\
GET & \url{/api/customer/history} & Customer & Lấy request, rental active, vé completed và request archive của khách. \\
POST & \url{/api/customer/rentals/:id/confirm-return} & Customer & Xác nhận bãi trả bằng tọa độ hiện tại trước khi staff xuất vé. \\
GET & \url{/api/gps-demo/users} & Public demo & Lấy danh sách tài khoản customer để \texttt{/gd} tìm và kéo GPS theo từng user. \\
GET & \url{/api/gps-demo/sessions} & Public demo & Lấy các phiên pending/active để \texttt{/gd} liên kết user với UC đang diễn. \\
GET / POST & \url{/api/gps-demo/position} & Public demo & Lưu/đọc tọa độ GPS demo theo từng customer giữa tab \texttt{/gd} và workspace khách hàng. \\
GET & \url{/api/staff/pending-requests} & Staff/Admin & Danh sách request chờ nhận xe. \\
POST & \url{/api/staff/handover} & Staff/Admin & Chuyển request thành rental sau khi xác nhận giấy tờ/cọc và kiểm tra GPS hiện tại của customer nằm trong bán kính bãi nhận. \\
GET & \url{/api/staff/active-rentals} & Staff/Admin & Danh sách rental đang chạy để xử lý trả xe. \\
POST & \url{/api/staff/return} & Staff/Admin & Trả xe sau xác nhận của customer, tính ticket, cập nhật vị trí/trạng thái xe. \\
GET & \url{/api/admin/bikes} & Staff/Admin & Tra cứu toàn bộ đội xe và trạng thái giữ chỗ. \\
GET & \url{/api/admin/reports} & Staff/Admin & Tổng hợp lượt thuê, doanh thu, phụ thu và fleet status. \\
GET & \url{/api/admin/bike-status-logs} & Staff/Admin & Đọc nhật ký đổi trạng thái xe, gồm mã xe, trạng thái cũ/mới, bãi, người đổi, lý do và thời gian. \\
POST & \url{/api/admin/users/:id/status} & Admin & Khóa/mở tài khoản để diễn account blocked alternative flow. \\
POST & \url{/api/admin/blocked-identities} & Admin & Thêm CCCD/CMND vào danh sách khóa vận hành. \\
DELETE & \url{/api/admin/blocked-identities/:identityNumber} & Admin & Gỡ CCCD/CMND khỏi danh sách khóa vận hành. \\
POST / PATCH & \url{/api/admin/stations} & Admin & Tạo/cập nhật bãi xe. \\
POST / PATCH & \url{/api/admin/bikes} & Admin & Tạo/cập nhật xe. \\
PATCH & \url{/api/admin/bikes/:id/status} & Staff/Admin & Đổi trạng thái xe kèm lý do và status log. \\
\bottomrule
\end{longtable}
\endgroup

Phiên đăng nhập được quản lý bằng session cookie riêng cho từng browser hoặc
browser context. Vì vậy khi trình bày, kịch bản thường có thể dùng một phiên
khách hàng, một phiên admin/operator và một tab \texttt{/gd} để kéo GPS người
dùng đang diễn và tua giờ hệ thống. Bộ smoke test vẫn kiểm tra mức tải demo tối đa gồm hai khách hàng,
hai tài khoản admin và một context \texttt{/gd}, bảo đảm các phiên không ghi đè
nhau. Seed demo cung cấp:
\begin{itemize}
  \item hai tài khoản khách: \texttt{customer@ecopark.test} và \texttt{resident@ecopark.test};
  \item hai tài khoản admin: \texttt{admin@ecopark.test} và \texttt{admin2@ecopark.test}.
\end{itemize}

Các quy tắc giá được triển khai đúng theo đề bài:
\begin{itemize}
  \item Giá cơ bản: 50k cho 1 giờ, 70k cho 2 giờ và 100k cho 3 giờ.
  \item Phụ thu trễ: 30k cho mỗi 30 phút phát sinh, làm tròn lên theo block 30 phút.
  \item Cư dân Ecopark hợp lệ được giảm 40\% trên phí thuê cơ bản.
  \item Tiền đặt cọc 200k và giấy tờ định danh được quản lý riêng, không tính là doanh thu vé thuê.
\end{itemize}

\section{Testing and evaluation}

\subsection{Local operation commands}
Repo cung cấp các lệnh vận hành chính:
\begin{itemize}
  \item \texttt{npm install}: cài dependency front-end/back-end.
  \item \texttt{npm run reset-db}: tạo lại SQLite database demo/local.
  \item \texttt{npm run dev}: chạy web app tại \texttt{http://127.0.0.1:4173}.
  \item \texttt{/gd}: mở demo director để kéo thả GPS người dùng theo tuyến đường demo và tua đồng hồ hệ thống.
  \item \texttt{npm run check}: kiểm tra cú pháp JavaScript cho app, server, test và script automation.
  \item \texttt{npm test}: chạy test nghiệp vụ back-end.
  \item \texttt{npm run test:coverage}: chạy test back-end kèm Node.js coverage report.
  \item \texttt{npm run smoke:ui}: kiểm tra render giao diện desktop/mobile bằng Playwright.
  \item \texttt{npm run smoke:uc}: chạy pipeline UI cho các luồng UC chính và một số alternative flow.
  \item \texttt{npm run screenshots:uc001}: sinh bộ ảnh UC001 và alternative flow 1 dùng trong report.
  \item \texttt{npm run docker:build}: kiểm tra container deployable cho target GCP.
\end{itemize}

\subsection{Test coverage}
\begin{longtable}{P{0.24\textwidth}P{0.36\textwidth}P{0.30\textwidth}}
\toprule
\textbf{Test layer} & \textbf{Coverage} & \textbf{Evidence/command} \\
\midrule
Backend API tests & Đăng ký, validation input, password policy, nâng cấp hash SHA-256 cũ sang PBKDF2, email verification, reset password, duplicate email/phone/CCCD, blocked identity, resident pending, stale resident-discount flag for visitor, held-bike availability, tạo request ngoài bán kính nhưng chặn handover theo GPS hiện tại, cancel request, exchange bike, staff station scope, geofence return, handover, capacity guard, return pricing, GPS demo theo từng customer, reports, audit logs, concurrent sessions. & \texttt{npm test} \\
Backend coverage & Coverage hiện tại theo Node.js built-in coverage: line 78.93\%, branch 69.31\%, functions 80.58\%. & \texttt{npm run test:coverage} \\
UI smoke tests & Auth/customer/admin trên 5 cỡ màn hình đại diện, GSAP load, Three.js canvas pixel, Leaflet marker, custom report dropdown, scale tiền của chart báo cáo, bảng vận hành rộng không scroll ngang, bảng vận hành hẹp dạng row-card, console error và horizontal overflow. & \texttt{npm run smoke:ui} \\
UC smoke tests & GPS context riêng, 2 customer + 2 admin sessions, chọn user trong \texttt{/gd} và kiểm marker customer map chạy theo, request/cancel/request again, staff handover, nhánh bấm xuất vé trước khi khách xác nhận để kiểm lý do bị chặn, customer confirm return, return ticket, chart dashboard, audit log và CSV report export. & \texttt{npm run smoke:uc} \\
Screenshot evidence & Bộ ảnh UC001 basic flow và alternative flow 1 được sinh từ chính app bằng Playwright. & \texttt{npm run screenshots:uc001} \\
CI/CD automation & GitHub Actions chạy static check, coverage, UI smoke, UC smoke và build Docker image deployable trên push/PR. & \texttt{.github/workflows/ci.yml} \\
Manual visual review & Screenshot auth, customer rental flow, operations workspace, mobile dock và GPS route. & Ảnh trong report UI section. \\
\bottomrule
\end{longtable}

Bộ test hiện kiểm tra đăng ký tài khoản, chống trùng email/phone/CCCD, chính
sách mật khẩu, nâng cấp hash cũ sang PBKDF2, xác minh email trước khi thuê,
reset mật khẩu demo-local, chặn CCCD/CMND bị khóa, thẻ cư dân pending thuê như
khách thường, loại xe đã giữ chỗ ra khỏi danh sách xe sẵn sàng, hủy yêu cầu
thuê và giải phóng xe, đổi xe trước handover, staff chỉ thấy/xử lý request đúng
bãi, chuyển yêu cầu thuê thành rental, chặn giảm sức chứa bãi dưới số xe hiện
có, tính vé trả xe có giảm giá cư dân/phụ thu trễ, sinh audit log trạng thái xe,
tổng hợp báo cáo và bốn phiên đăng nhập độc lập cho hai khách hàng cùng hai
admin.
Smoke test giao diện bằng Playwright mở các trạng thái desktop/mobile, kiểm tra
GSAP vendor đã tải, canvas Three.js có render, bản đồ Leaflet được mount với marker
thật, không phát sinh console error và không có tràn ngang ở cấp trang.
Smoke test UC mở demo director \texttt{/gd} ở context riêng, kiểm tra marker người dùng và trạng
thái đủ gần bãi trong kịch bản tải tối đa gồm bốn browser context cho hai khách
hàng và hai admin, đăng nhập bằng tài khoản cư dân, kiểm tra marker vị trí người
dùng, thử vị trí nhập tay ngoài phạm vi, gửi và hủy yêu cầu thuê, gửi lại yêu
cầu, đăng nhập Staff để giao xe và trả xe tại bãi được phân công, xuất ticket,
kiểm tra dashboard chart/audit log và tải báo cáo CSV; luồng trả khác bãi vẫn
được backend/UI hỗ trợ khi Admin/Operator xử lý toàn hệ thống.

\subsection{Evaluation and remaining limitations}
\begin{itemize}
  \item Sản phẩm đáp ứng đầy đủ 5 use case chính và các alternative flow quan
  trọng cần demo: GPS lỗi/nhập tay, hết xe trong phạm vi, hủy request, sai cọc/
  giấy tờ, trả khác bãi, xe hỏng và lọc báo cáo.
  \item Back-end đủ dùng cho demo/nộp môn học, đã dùng PBKDF2 salted password
  hash và mã xác minh/reset lưu dạng hash. Nếu chuyển production cần tích hợp
  dịch vụ gửi email thật, persistent session store, migration versioned và
  backup database.
  \item Dashboard vận hành đã có metric, chart doanh thu theo kỳ, phân bổ trạng
  thái đội xe, tổng thu/phụ thu cùng scale tiền, bảng chi tiết, CSV theo bộ lọc và audit timeline
  cho thay đổi trạng thái xe.
  \item Route \texttt{/gd} là công cụ thuyết trình để diễn vị trí người dùng theo đường
  và tua giờ nội bộ, không phải tracking GPS thật của thiết bị gắn trên xe hoặc điện thoại.
\end{itemize}
